\section{Objetivos}
\label{sec:objetivos}

El objetivo principal de este Trabajo Fin de Grado (TFG) es el desarrollo de una aplicación web, denominada \texttt{MoodNest}, para el seguimiento diario y análisis del estado de ánimo personal. La finalidad es proporcionar al usuario una herramienta para el autoconocimiento, permitiéndole registrar sus emociones y relacionarlas con sus actividades diarias.

Para ello, el proyecto debe cumplir otros objetivos más específicos como:
\begin{itemize}
    \item Fomentar la constancia del usuario. Para ello se ofrecerá una interfaz simple, intuitiva y, sobre todo, rápida, que minimice el esfuerzo del registro diario.
    
    \item Garantizar la privacidad y gestión de usuarios. Cualquier persona podrá registrarse de forma segura mediante un usuario y contraseña, asegurando que sus datos emocionales permanezcan privados y accesibles solo para él.
    
    \item Capacidad para gestionar los registros de ánimo. El sistema permitirá al usuario crear, consultar, editar y eliminar sus entradas diarias de estado de ánimo.
    
    \item Permitir la personalización de la escala de ánimo. Se debe permitir al usuario definir sus propios estados de ánimo (por ejemplo, nombres, colores y un valor numérico asociado a una escala, también ajustable).
    
    \item Implementar un sistema de etiquetado (\textit{tags}) gestionado en su totalidad por el usuario, pudiendo crear sus propias etiquetas (ej. \texttt{\#Trabajo}, \texttt{\#Dormir}, \texttt{\#Social}) y personalizarlas.

    \item Desarrollar un módulo de análisis y visualización de datos preciso. Para ello, la aplicación procesará el histórico de registros para identificar correlaciones entre las actividades diarias (\textit{tags}) y el estado de ánimo del usuario. El sistema generará gráficos de evolución y distribución que revelen patrones fiables sobre el impacto positivo o negativo de estos factores contextuales en el bienestar global.    
    
    \item Garantizar la consulta histórica de los datos. El usuario deberá poder navegar por sus registros pasados de manera accesible, por ejemplo, a través de una vista de calendario o un listado filtrable.
\end{itemize}